\section{Introduction}\label{sec:introduction}

Pricing American-style options requires implied volatility (IV) as an input, yet IV itself is an empirical quantity extracted from observed option prices~\cite{black1973}. This circularity is a practical obstacle for forward pricing and scenario simulation: existing approaches either take IV from the market and feed it into a pricing model at every step of a simulated path, or generate price paths without a consistent IV surface to price contracts against. Parametric IV models such as SABR~\cite{hagan2002} and the SVI parameterization~\cite{gatheral2004} treat the IV surface as an object to be fitted to market quotes, not generated from first principles; they produce smooth interpolations of observed smiles but cannot generate new IV dynamics for unobserved market conditions. Local volatility models~\cite{dupire1994,derman1994} recover a deterministic volatility surface consistent with a single cross-section of option prices, yet the resulting surface is static and cannot produce the stochastic IV fluctuations observed in real markets. Stochastic volatility models such as that of \citet{heston1993} generate realistic variance paths, but calibrating them requires observed option prices or IV as input, reintroducing the very circularity that a forward pricing pipeline must avoid. The result is a methodological gap: no existing computational pricer separates one-time calibration of an IV-shape function on a historical corpus from forward simulation that propagates a self-consistent IV surface path-conditionally, without ingesting new market-IV inputs along the simulated path.

This gap has direct consequences for forward pricing and scenario simulation. Empirical studies of the IV surface~\cite{cont2002} have established that smile, skew, and term-structure effects interact with market regimes, so a forward simulator has to reproduce that structure dynamically rather than freeze a single cross-section. Calibrating Heston or SABR to a market snapshot~\cite{hagan2002,gatheral2004} produces a self-consistent surface at one instant but locks the parameters to that snapshot, so forward-simulated option prices drift from the structural return dynamics whenever the underlying regime moves. A multi-date calibration of a shared shape function on a historical ladder corpus, by contrast, fixes the calibration cost up front and lets the forward simulator propagate IV without re-ingesting market quotes at each step. Historical resampling~\cite{glasserman2004} preserves realized distributional features but cannot produce out-of-sample tail scenarios consistent with a structural model of returns, and generative neural approaches~\cite{wiese2020} reproduce equity stylized facts without producing self-consistent option-price and IV pairs, so any downstream pricing step has to reintroduce the very calibration loop the simulator was meant to avoid. A practitioner who wants path-conditional American option prices, finite-difference Greeks, and forward IV from one coherent simulation therefore has to either accept the static-surface compromise or stitch together inconsistent models~\cite{buehler2019}, neither of which serves scenario analysis or backtesting pipelines that depend on option-level outputs.

In this study, we developed a computational pricing pipeline in which IV emerged as a path-conditional \emph{output} of forward simulation, after the IV-shape function was calibrated once on a historical option-ladder corpus and held fixed thereafter. The pipeline coupled a Jump Hidden Markov Model (JumpHMM)~\cite{jumphmm2025} that produced multi-asset price paths with realistic heavy tails, negligible linear autocorrelation, and persistent volatility clustering~\cite{cont2001}, with cross-asset dependence imposed via a Student-$t$ copula~\cite{demarta2005}; a modified Heston stochastic variance process~\cite{heston1993} that converted those price paths into IV paths, where the mean-reversion target $\theta$ was a function of the HMM regime state, days to expiration (DTE), moneyness, and an aggregate market-mood indicator; and a Cox-Ross-Rubinstein (CRR) binomial tree~\cite{cox1979} that converted the pre-computed IV into American option prices with early exercise. The central innovation was a hybrid $\theta$-function that linked the Heston mean-reversion target to the JumpHMM regime state, producing regime-dependent volatility dynamics. An equilibrium initialization $v_0 = \theta(t\!=\!0)$ gave each (ticker, strike, DTE) triple its own initial IV; once the shape function was calibrated against the historical ladder corpus, the forward simulator propagated smile, skew, and term-structure dynamics path-conditionally without re-ingesting market IV along the simulated path. Within this pipeline, the shape function that controlled the smile and term structure of $\theta$ admitted a family of representations, and the modeling choice at each level of the family traded off model capacity for sharing across tickers. We therefore calibrated the shape function through a hierarchy of increasingly expressive forms: a parsimonious five-parameter functional form on single-ticker and cross-ticker semiconductor option chains, a globally shared neural surrogate that replaced the analytic form with a data-driven representation, and a sector-specific neural surrogate that let within-sector shapes specialize. The same decomposition supported a further relaxation to per-ticker networks once per-name sample sizes grew, without changing the downstream pipeline. We demonstrated this progression on a 31-ticker, fifteen-date, six-sector option ladder and on an eight-day temporal-holdout subset of the same capture, the latter of which isolated scheduled earnings events as the dominant source of the temporal generalization gap and motivated calendar-derived earnings-distance and same-sector peer-distance features that recovered a portion of that gap. The framework also produced realistic forward IV paths and American option prices with an endogenous crash-protection premium. The pipeline was implemented as an open-source Julia package so that the calibration, forward simulation, and pricing steps could be reproduced and reused end to end.
